\documentclass[tikz,border=6pt]{standalone}
\usepackage{ctex}
\usepackage{amsmath}
\usepackage{amssymb}
\usepackage{pgfplots}
\pgfplotsset{compat=1.18}
\usetikzlibrary{calc,angles,quotes,arrows.meta,positioning,decorations.markings}
\begin{document}
\begin{tikzpicture}
\begin{axis}[
    width=13cm, height=9cm,
    xlabel={$x$}, ylabel={$y$},
    xmin=-3, xmax=2, ymin=-0.4, ymax=7.6,
    xtick={-3,-2,-1,0,1,2},
    ytick={0,1,2,3,4,5,6,7},
    axis lines=middle,
    grid=major, grid style={gray!18},
    tick label style={font=\small},
    label style={font=\small},
    legend style={font=\small, draw=gray!40, at={(0.4,0.97)}, anchor=north west},
    enlarge x limits=0.02,
    restrict y to domain=-1:9,
]
% 目标曲线 e^x（粗黑虚线）
\addplot[black, very thick, dashed, forget plot, domain=-3:2, samples=140] {exp(x)};
% n=2: (1+x/2)^2
\addplot[orange!90!black, thick, forget plot, domain=-3:2, samples=140] {(1+x/2)^2};
% n=5: (1+x/5)^5
\addplot[blue!75!black, thick, forget plot, domain=-3:2, samples=160] {(1+x/5)^5};
% n=1: 1+x （直线，最后画并用点状线，确保从其他曲线上方透出可见）
\addplot[red!85!black, very thick, densely dotted, forget plot, domain=-3:2, samples=2] {1+x};
% 图例：手工按 n 递增排列（与绘制顺序解耦，确保颜色与标签一一对应）
\addlegendimage{black, very thick, dashed}\addlegendentry{$y=e^x$（极限曲线）}
\addlegendimage{red!85!black, very thick, densely dotted}\addlegendentry{$n=1:\ (1+x)$}
\addlegendimage{orange!90!black, thick}\addlegendentry{$n=2:\ (1+\tfrac{x}{2})^2$}
\addlegendimage{blue!75!black, thick}\addlegendentry{$n=5:\ (1+\tfrac{x}{5})^5$}
% 注释框：随 n 增大逐步逼近 e^x（置于左下空白区，避免被右侧曲线挤裁）
\node[anchor=north west, align=left, font=\small,
      fill=white, fill opacity=0.85, text opacity=1, inner sep=3pt, draw=gray!30]
    at (axis cs:-2.92,7.3)
    {$\Big(1+\dfrac{x}{n}\Big)^{n}\to e^{x}$\\[2pt]$n$ 越大，多项式曲线\\越贴近 $e^x$（离散$\to$连续）};
% 共同点 (0,1)
\addplot[only marks, mark=*, mark size=1.8pt, black] coordinates {(0,1)};
\node[anchor=north east, font=\footnotesize] at (axis cs:0,1) {$(0,1)$};
\end{axis}
\end{tikzpicture}
\end{document}
